%
% agradecimentos.tex
% Capítulo não numerado de Agradecimentos
%
\chapter*{Agradecimentos}
\addcontentsline{toc}{chapter}{Agradecimentos}
\label{cha:agradecimentos}
\enlargethispage{2.5cm}
{\setlength{\parskip}{3pt}

A realização deste trabalho não teria sido possível sem o contributo, directo ou indirecto, de um conjunto de pessoas a quem devo reconhecimento genuíno.

Ao \textbf{Professor Nuno Datia} e ao \textbf{Professor Michel Vorenhout}, orientadores deste projecto, expresso a minha gratidão pela disponibilidade, pelas observações e pela liberdade intelectual que me concederam para explorar caminhos pouco convencionais. A qualidade deste trabalho deve muito à exigência que colocaram e à confiança que depositaram. Ao \textbf{Instituto Superior de Engenharia de Lisboa} e ao \textbf{Instituto Politécnico de Lisboa}, pela formação académica e pelo enquadramento institucional que tornaram possível um projecto com esta dimensão e ambição.

Aos colegas com quem partilhei percurso ao longo destes anos, agradeço as conversas, os debates técnicos e a camaradagem que tornam o processo de aprender menos solitário e mais fértil. A disponibilidade para discutir ideias, questionar abordagens e partilhar dificuldades foi, muitas vezes, o que impediu que um obstáculo se tornasse um impasse.

Aos amigos de longa data, o reconhecimento pela paciência e pela presença constante. A um projecto desta natureza não se exige apenas competência técnica --- exige-se também resistência, e essa constrói-se com o apoio de quem se mantém próximo independentemente das circunstâncias.

À minha família, o agradecimento mais profundo e mais difícil de circunscrever em palavras. A dedicação exigida por um trabalho desta envergadura impõe custos que recaem, inevitavelmente, sobre aqueles que estão mais próximos. A paciência, o suporte incondicional e a crença no valor daquilo que estava a ser construído foram indispensáveis.

\vspace{0.5cm}

\noindent A todos os que contribuíram para este percurso:

\vspace{0.2cm}

\begin{center}
\begin{minipage}{0.85\textwidth}
\centering
\textit{Família:} Ilesh Gamanbhai, Jignaben Mistry, Divyesh Govind.

\vspace{4pt}
\textit{Orientação e corpo docente:} Professor Nuno Datia, Professor António Serrador, Professor Michel Vorenhout.

\vspace{4pt}
\textit{Colegas e amigos:} Pedro Gomes, Guilherme Coutinho, Martim Monteiro, Alexandre Tavares, André Nunes, Francisco Tavares, Rodrigo Sampaio, Simão Afonso, Fábio Nibau, Maria Lopes, Júlia Santana, Vasco Garcia.
\end{minipage}
\end{center}

\nopagebreak[4]
\vspace{0.4cm}
\begin{flushright}
Radesh Ilesh Gamanbhai Govind
\end{flushright}

}
